% sections/en/3-problem.tex -- problem statement, assumptions, notation, exact model, fluid relaxation, gap identity.
\section{Problem statement and fluid--quota geometry}\label{sec:problem}

\subsection{Problem description}\label{sec:problem-desc}
A shop has a machine set $M$ and a set $\mathcal{R}$ of job types. Type $r$ consists of $N_r \ge 1$ identical copies,
each processed along a fixed routing of $K_r$ operations. The type-operation $(r,k)$ can be processed on any machine of
its candidate set $M_{rk} \subseteq M$ with processing time $p_{rki}$; all copies of a type share candidate sets and
processing times, but each copy must satisfy the precedence constraints of its own routing. The decisions are the machine
assignment of every copy-operation and the processing order on every machine. All data are known in advance and the
schedule is produced once, so the setting is static and deterministic. When $N_r = 1$ for every $r$ the problem reduces to
the standard flexible job-shop scheduling problem (FJSP).

The following assumptions are made. A1: processing times are deterministic and known (standard). A2: a machine processes
at most one operation at a time (standard). A3: processing is non-preemptive (standard). A4: all copies are available at
time zero and there are no dynamic arrivals (standard; stated under limitations). A5: setup times, transportation and lot
splitting are ignored (simplification of this paper; stated under limitations).

\subsection{Notation}\label{sec:notation}
Table~\ref{tab:notation} lists the symbols; a symbol keeps one meaning throughout the paper, and the appendix (Table~\ref{tab:glossary}) collects the terminology. A copy-operation is
$o = (r, c, k)$ with copy index $c$, and $p_{oi} = p_{rki}$. The shortest route length of a type is
$\ell_r^{\min} = \sum_k \min_{i \in M_{rk}} p_{rki}$ and $B$ is a sufficiently large constant.

\begin{table}
\centering\small
\caption{Notation.}\label{tab:notation}
\begin{tabularx}{\linewidth}{@{}>{\hsize=0.72\hsize}L>{\hsize=1.28\hsize}L@{}}
\toprule
\textbf{Symbol} & \textbf{Meaning} \\
\midrule
$\mathcal{R},\,M$ & job-type set, machine set ($R = \lvert\mathcal{R}\rvert$, $m = \lvert M\rvert$) \\
$r,\,c,\,i$; $o,\,o'$ & type, copy and machine indices; copy-operations $o = (r,c,k)$ \\
$N_r,\,K_r$ & copies of type $r$; operations per copy \\
$M_{rk},\,p_{rki}$ & candidate machines and processing time of $(r,k)$ \\
$\ell_r^{\min},\,B$ & shortest route length of type $r$; big-M constant \\
$\bar N,\,V_N,\,F,\,V_p$ & mean multiplicity, multiplicity dispersion, flexibility, processing-time heterogeneity \\
$x_{oi},\,s_o,\,y_{oo'}$ & assignment, start time and sequencing variables \\
$C_{\max}$ & makespan \\
$q_{rki},\,L^F,\,\lambda_i$ & fluid variable, fluid lower bound, machine-load dual (feature only) \\
$n_{rki},\,\kappa_{rk},\,\chi$ & integer quota, priority key, individual $\chi = (\mathbf{n}, \boldsymbol{\kappa})$ \\
$L^Q(\mathbf{n})$ & integer-quota bound \\
$G_Q,\,G_S$; $g_Q,\,g_S$ & rounding loss, sequencing residual; their $L^F$-normalized values \\
$B_i,\,\bar B$ & rounding bound of machine $i$ and its maximum \\
$\mathcal{A},\,a,\,T_a$ & candidate pool, action, action transform \\
$\hat G_S(a)$ & predicted sequencing residual of action $a$ \\
$K_L,\,K_H,\,K_R$ & candidates taken from the learner, the fluid heuristic and at random \\
\bottomrule
\end{tabularx}
\end{table}

\subsection{Exact model}\label{sec:model}
The job-indexed model is
\begin{align}
\min\quad & C_{\max} \label{eq:obj}\\
\text{s.t.}\quad & \textstyle\sum_{i \in M_{rk}} x_{oi} = 1, && \forall\, o \label{eq:assign}\\
& s_{(r,c,k+1)} \ge s_o + \textstyle\sum_{i} p_{oi} x_{oi}, && \forall\, o:\, k < K_r \label{eq:prec}\\
& s_{o'} \ge s_o + p_{oi} - B(3{-}x_{oi}{-}x_{o'i}{-}y_{oo'}), && \forall\,i,\,o{\ne}o' \label{eq:disj}\\
& C_{\max} \ge s_o + \textstyle\sum_{i} p_{oi} x_{oi}, && \forall\, o:\, k = K_r \label{eq:cmax}\\
& x_{oi},\, y_{oo'} \in \{0,1\},\quad s_o \ge 0. \label{eq:dom}
\end{align}
Constraint~\eqref{eq:assign} assigns every copy-operation to exactly one candidate machine, \eqref{eq:prec} enforces the
routing of each copy, \eqref{eq:disj} is the big-M disjunction on a shared machine ($y_{oo'} = 1$ when $o$ precedes $o'$),
\eqref{eq:cmax} defines the makespan and \eqref{eq:dom} the domains. The model is complete enough to be handed to CP-SAT
\cite{ortools}, which is how the exact reference solutions of the small tier are produced.

\subsection{Quota-type fluid relaxation}\label{sec:fluid}
The relaxation keeps the per-type totals and drops the discreteness of the copies:
\begin{align}
L^F = \min\quad & L \label{eq:fluid-obj}\\
\text{s.t.}\quad & \textstyle\sum_{i \in M_{rk}} q_{rki} = N_r, && \forall\, r, k \label{eq:fluid-cons}\\
& \textstyle\sum_{r}\sum_{k:\, i \in M_{rk}} p_{rki}\, q_{rki} \le L, && \forall\, i \label{eq:fluid-load}\\
& L \ge \ell_r^{\min},\quad q_{rki} \ge 0, && \forall\, r, k, i. \label{eq:fluid-dom}
\end{align}
Constraint~\eqref{eq:fluid-cons} routes the $N_r$ copies of type-operation $(r,k)$ to its candidate machines as a
continuous mass, \eqref{eq:fluid-load} bounds the fluid load of every machine by the cycle length and
\eqref{eq:fluid-dom} keeps the cycle length above every shortest route. $L^F$ is the fluid lower bound
(Proposition~\ref{prop:fluidlb}) and $q^*_{rki}/N_r$ is the macroscopic flow ratio of $(r,k)$. The machine-load duals
$\lambda_i$ are used only as ranker features: at the fluid optimum they are close to uniform and carry little guidance.

\subsection{Integer quotas and the gap identity}\label{sec:gap}
An integer quota $\mathbf{n}$ assigns $n_{rki}$ copies of $(r,k)$ to machine $i$, with $n_{rki}$ a non-negative integer,
$n_{rki} = 0$ for $i \notin M_{rk}$ and $\sum_i n_{rki} = N_r$. It induces the quota bound
\begin{align}
L^Q(\mathbf{n}) &= \max\{ \max_i \textstyle\sum_{r,k:\, i \in M_{rk}} p_{rki} n_{rki},\, \max_r \ell_r^{\min} \}, \label{eq:LQ}\\
C_{\max}(\chi) - L^F &= \underbrace{L^Q(\mathbf{n}) - L^F}_{G_Q} + \underbrace{C_{\max} - L^Q(\mathbf{n})}_{G_S}. \label{eq:decomp}
\end{align}
Identity~\eqref{eq:decomp} is the anchor of the paper: for any feasible individual the fluid gap splits exactly into the
rounding loss $G_Q$, which depends only on the quota, and the sequencing residual $G_S$. The normalized quantities
$g_Q = G_Q/L^F$ and $g_S = G_S/L^F$ serve for cross-instance comparison and as ranker labels. Two properties hold without
any assumption on the search.

\begin{proposition}[Fluid lower bound]\label{prop:fluidlb}
Every schedule satisfying \eqref{eq:assign}--\eqref{eq:dom} has $C_{\max} \ge L^F$.
\end{proposition}
\begin{proof}
Let $q_{rki} = \sum_c x_{(r,c,k),i}$. By \eqref{eq:assign}, $\sum_i q_{rki} = N_r$, which is \eqref{eq:fluid-cons}. The
operations assigned to machine $i$ do not overlap by \eqref{eq:disj}, so their total length $\sum_{r,k} p_{rki} q_{rki}$ is at
most $C_{\max}$, which is \eqref{eq:fluid-load} with $L = C_{\max}$. Each copy completes its $K_r$ operations in series by
\eqref{eq:prec}, so $C_{\max} \ge \ell_r^{\min}$ and \eqref{eq:fluid-dom} holds. Hence $(q, C_{\max})$ is feasible for
\eqref{eq:fluid-obj}--\eqref{eq:fluid-dom} and $L^F \le C_{\max}$.
\end{proof}

\begin{proposition}[Non-negative decomposition]\label{prop:decomp}
For every feasible schedule decoded from a quota individual $\chi = (\mathbf{n}, \boldsymbol{\kappa})$, identity
\eqref{eq:decomp} holds with $G_Q(\chi) \ge 0$ and $G_S(\chi) \ge 0$.
\end{proposition}
\begin{proof}
The identity is immediate. With $q = \mathbf{n}$ and $L = L^Q(\mathbf{n})$, constraint \eqref{eq:fluid-cons} holds by
the quota totals, \eqref{eq:fluid-load} by the first term of \eqref{eq:LQ} and \eqref{eq:fluid-dom} by the second, so
$L^Q(\mathbf{n}) \ge L^F$ and $G_Q \ge 0$. After decoding, the operations on machine $i$ do not overlap, so
$\sum_{r,k} p_{rki} n_{rki} \le C_{\max}$ for every $i$, and every copy runs its route in series, so
$C_{\max} \ge \ell_r^{\min}$; together $C_{\max} \ge L^Q(\mathbf{n})$ and $G_S \ge 0$.
\end{proof}
